
\documentclass{../concepts/prepareprotokol} %všechny balíčky pro protokol + titulní strana
\usepackage{graphicx}
\usepackage{booktabs}
\usepackage{amsmath}
\usepackage{newfloat}
\usepackage{caption}
\usepackage{siunitx}
\sisetup{inter-unit-product =$\cdot$}
\sisetup{separate-uncertainty=true}

\praktikum{II} %číslo praktika (I, II nebo III)
\autor{David Němec} %jméno autora
\datum{10.\,11.\,2025} %datum měření
\cislo{4} %číslo úlohy
\nazev{Měření malých odporů} %název úlohy

\DeclareFloatingEnvironment[fileext=los, placement={htbp}, name=Obrázek]{schema}

\newcommand{\nobars}{
Chybové úsečky nebyly pro svou malou velikost pro přehlednost vykresleny.}

\newcommand{\relunc}[1]{
\left( \frac{\sigma_{#1}}{#1} \right)^2
}


\begin{document}
\renewcommand{\figurename}{Graf}
\renewcommand{\refname}{Literatura}

\maketitle

% ========================== PRACOVNÍ ÚKOL =========================
\section{Pracovní úkol}
\begin{enumerate}
    \item Změřte průměry a délky drátů na pracovní desce.

    \item Změřte odpory ve čtyřbodovém zapojení pomocí multimetru KEITHLEY 2010.
    
    \item Změřte odpor drátů Wheatstoneovým a Thomsonovým můstkem Metra - MTW. Vysvětlete rozdíly ve výsledcích měření. Současně určete odpor přívodních vodičů a odpor na svorkách v případě měření Wheatstoneovým můstkem.
    
    \item Určete měrný odpor jednotlivých vzorků i s příslušnou chybou výsledku. Stanovené hodnoty porovnejte s hodnotami uváděnými v tabulkách.

\end{enumerate}

% ============================= TEORIE ===============================
\section{Teorie}

\subsection{Čtyřbodové zapojení rezistoru}

Pro přesnější měření odporu rezistoru lze využít čtyřbodového zapojení podle schématu na obrázku \ref{sch:four}. Pokud je voltmetr připojen k napěťovým svorkám $b$ a $b'$, nebude měření ovlivněno spádem napětí na přívodních vodičích mezi rezistorem a proudovými svorkami $a$ a $a'$. Za předpokladu nízkých hodnot $R$ oproti odporu voltmetru poteče přívodními vodiči k napěťovým svorkám zanedbatelný proudu a ani odpor těchto vodičů se neuplatní \cite{studijni}.

\subsection{Wheastonův můstek}

Nízké hodnoty odporu studovaného rezistoru lze měřit pomocí Wheastonova můstku. Jeho schéma zapojení je na obrázku \ref{sch:wheast}. Je-li můstek v rovnováze, tj. galvanometrem G neteče proud, platí pro poměry jednotlivých odporů $R_1$, $R_2$, $R_3$ a $R_4$ vztah \cite{studijni}
\begin{equation}
    \frac{R_1}{R_2} = \frac{R_3}{R_4} \, .
\end{equation}
Zafixováním poměru $R_3$\,/\,$R_4$ na vysoké hodnotě a postupným měněním odporu $R_2$ tak, aby výchylka galvanometru byla nulová, lze pak snadno změřit zkoumaný nízký odpor $R_1$. V jeho hodnotě však bude v důsledku tohoto zapojení započten i odpor přívodních drátů.

\subsection{Thomsonův můstek}

Přesnější měření lze provést pomocí Thomsonova můstku, jehož schéma zapojení je na obrázku \ref{sch:four}. Pokud bude opět můstek v rovnováze (galvanometrem G nepoteče proud) a pro odpory $P$, $Q$, $p$ a $q$ bude platit
\begin{equation}
    \frac{p}{P} = \frac{q}{Q} \, ,
\end{equation}
lze zkoumaný odpor $R_X$ určit jednoduše jako \cite{studijni}
\begin{equation} \label{eq:thomson}
    R_X = \frac{Q}{P} R_N \, .
\end{equation}
Při čtyřbodovém zapojení odporů $R_N$ a $R_X$ již nebude měření ovlivněno odporem jejich přívodních drátů. Za předpokladu nízkých hodnot $R_N$ a $R_X$ oproti odporům $P$, $Q$, $p$ a $q$, budou také nízké proudy $I$ a $i$ oproti proudu $I_0$ a odpory přívodních vodičů k můstku se díky tomu také v měření neuplatní \cite{studijni}.

\subsection{Měrný odpor}

Měrný odpor $\rho$ drátu průřezu $S$ a délky $l$ je definován vztahem \cite{studijni}
\begin{equation} \label{eq:rho}
    \rho = \frac{R S}{l} \, ,
\end{equation}
kde $R$ je celkový odpor drátu. Je-li příčný řez drátu kruhový, lze průřez $S$ vypočítat jako
\begin{equation} \label{eq:s}
    S = \frac{\pi D^2}{4} \, ,
\end{equation}
kde $D$ je průměr drátu.


\subsection{Metoda měření}

Nejprve byly zjištěny rozměry zkoumaných drátů - průměr a délka, přičemž délka byla měřena mezi 
napěťovými kontakty na obou koncích drátu. Pro větší přesnost byl průměr měřen na několika místech drátu.

Při měření multimetrem KEITHLEY 2010 bylo použito čtyřbodové zapojení rezistoru (obr. \ref{sch:four}), 
napěťové a proudové kontakty byly propojeny se vstupními svorkami multimetru podle pokynů v \cite{pokyny}.

Při měření odporu pomocí Wheastonova a Thomsonova můstku byly nastaveny hodnoty zbylých odporů 
($R_3$ a $R_4$ pro Wheastonův můstek a $P$ pro Thomsonův můstek) tak, aby hodnota na odporové 
dekádě v jednotkách \si{\ohm} odpovídala hodnotě zkoumaného odporu v jednotkách \si{\milli \ohm} 
(viz \ref{wheaston} a \ref{thomson}). Jako první odhad optimální hodnoty odporu na odporové dekádě 
posloužila hodnota odporu určená pomocí multimetru KEITHLEY 2010. Otáčením postupně jemnějších koleček 
dekády byla vyladěna výsledná hodnota tak, aby byl můstek v rovnováze, tedy aby galvanometrem netekl 
proud a jeho výchylka byla nulová.

Při měření s Wheastonovým můstkem byl galvanometr do obvodu zapojen vždy jen na krátkou dobu, aby proud
obvodem netekl dlouho a jednotlivé rezistory (a dráty) se nezahřívaly a něměnily tak svůj odpor.

\begin{schema}[htbp]
    \centering
    \includegraphics[width=0.5\linewidth]{uloha4/fourpoints.png}
    \caption{Čtyřbodové zapojení rezistoru \cite{studijni}}
    \label{sch:four}
\end{schema}

\begin{schema}[htbp]
    \centering
    \includegraphics[width=0.5\linewidth]{uloha4/wheaston.png}
    \caption{Wheastonův můstek \cite{studijni}}
    \label{sch:wheast}
\end{schema}

\begin{schema}[htbp]
    \centering
    \includegraphics[width=0.5\linewidth]{uloha4/thomson.png}
    \caption{Thomsonův můstek \cite{studijni}}
    \label{sch:thomson}
\end{schema}


\subsection{Měřicí přístroje a jejich nejistoty} \label{pomucky}
\begin{enumerate}
    \item Mikrometr s nejmenším dílkem 0.01 mm:
    
    Mikrometrem byly měřeny průměry jednotlivých drátů z různých materiálů. Nejistotu naměřené hodnoty odhaduji jako velikost nejmenšího dílku, tedy 0.01 mm.

    \item Pásový metr s nejmenším dílkem 1 mm:

    Metrem byly měřeny délky jednotlivých drátů. Nejistotu naměřené hodnoty odhaduji jako velikost nejmenšího dílku, tedy 1 mm. Některé dráty však nebyly mezi svorkami napnuty úplně dokonale, proto je nejistota jejich délky větší (viz \ref{rozmery}).

    \item Multimetr KEITHLEY 2010:

    Multimetrem se měřily odpory drátů ve čtyřbodovém zapojení. Z manuálu k použití přístroje v příloze v \cite{pokyny} je nejistota měření odporu (pro rozsah \SI{10}{\ohm}, na kterém bylo měřeno) 60 ppm z naměřené hodnoty + 9 ppm z rozsahu.

    \item Odporová dekáda:

    Nastavováním odporu na odporové dekádě byla hledána rovnováha na Wheastonově a Thomsonově můstku. Nejistota 
    nastaveného odporu byla odhadnuta jako o řád větší než největší řád, který již neměl vliv na výchylku galvanometru 
    (tj. výchylka byla nulová a nijak se něměnila ani při otáčení příslušného kolečka odporové dekády, např. pokud se otáčením
    kolečka odpovídajícího jednotkám \si{\ohm} výchylka neměnila, byla nejistota odhadnuta jako \SI{10}{\ohm}). Nejistota
    dekády uvedená na její přední straně se ukázala být zanedbatelná vůči takto odhadnuté nejistotě.

    \item Termohygrobarometr Commeter C4130:

    termohygrobarometrem byla zjištěna teplota okolí během měření, aby experimentálně určené hodnoty měrného odporu mohly být srovnány s tabulkovými. Nejistota teploty byla určena z manuálu přístroje jako \SI{0.4}{\celsius}.

\end{enumerate}

% ======================== VÝSLEDKY MĚŘENÍ ===========================
\section{Výsledky měření}

\subsection{Rozměry drátů a teplota okolí} \label{rozmery}

Průměry jednotlivých drátů byly změřeny mikrometrem na deseti místech a z nich byl 
vypočtena střední hodnota (aritmetický průměr). Naměřené hodnoty a jejich střední 
hodnoty shrnuje tabulka \ref{tab:prumer}. Nejistota středních hodnot byla vypočtena jako 
\begin{equation}
    \sigma_{\bar{D}} = \sqrt{\frac{1}{9} \sum_{i=1}^{10} (\bar{D} - D_i)^2 + \sigma_{mic}^2} \,
\end{equation}
kde $\bar{D}$ je střední hodnota průměru drátu, $D_i$ je i-tá hodnota průměru a 
$\sigma_{mic}=\SI{0.01}{\milli\metre}$ je nejistota měřidla.

Vypočtené střední hodnoty průměrů drátů a jejich délky ukazuje tabulka \ref{tab:delky}. 
Nejistota určení délky byla u~většiny drátů vzata jako nejistota pásového metru. Měděný 
drát však zjevně nebyl mezi svorkami dokonale napnutý, proto jeho délka i její nejistota 
byly odhadnuty jako o něco větší, než byla přímá vzdálenost mezi svorkami.

\begin{table}[!hbtp]
\centering
\caption{Změřené průměry drátů}
\begin{tabular}{c|cccccc}
 \toprule
 materiál & wolfram & měď & kanthal & železo & mosaz & chromnikl \\
 \midrule

 1 & $0.69$ & $1.34$ & $0.5$ & $0.41$ & $0.59$ & $1.02$ \\
 2 & $0.68$ & $1.35$ & $0.5$ & $0.41$ & $0.6$ & $1.01$ \\
 3 & $0.69$ & $1.35$ & $0.5$ & $0.41$ & $0.6$ & $0.99$ \\
 4 & $0.69$ & $1.34$ & $0.49$ & $0.41$ & $0.6$ & $1.0$ \\
 5 & $0.69$ & $1.35$ & $0.5$ & $0.41$ & $0.59$ & $1.01$ \\
 6 & $0.68$ & $1.35$ & $0.5$ & $0.41$ & $0.61$ & $1.0$ \\
 7 & $0.69$ & $1.35$ & $0.51$ & $0.41$ & $0.59$ & $1.0$ \\
 8 & $0.69$ & $1.35$ & $0.51$ & $0.41$ & $0.6$ & $0.99$ \\
 9 & $0.68$ & $1.34$ & $0.5$ & $0.41$ & $0.6$ & $1.0$ \\
 10 & $0.69$ & $1.34$ & $0.5$ & $0.41$ & $0.59$ & $1.02$ \\
 \midrule
 $\bar{D}$ (mm) & $0.687 \pm 0.011$ & $1.346 \pm 0.011$ & $0.501 \pm 0.011$ & $0.410 \pm 0.010$ & $0.597 \pm 0.012$ & $1.004 \pm 0.015$ \\
 
 \bottomrule
\end{tabular}
\label{tab:prumer}
\end{table}

\begin{table}[!hbtp]
\centering
\caption{Změřené průměry a délky drátů}
\begin{tabular}{c|cccccc}
 \toprule
 materiál & wolfram & měď & kanthal & železo & mosaz & chromnikl \\
 \midrule
 D (\si{\milli\metre}) & $0.687 \pm 0.011$ & $1.346 \pm 0.011$ & $0.501 \pm 0.011$ & $0.410 \pm 0.010$ & $0.597 \pm 0.012$ & $1.004 \pm 0.015$ \\
 l (\si{\centi\metre}) & $89.9 \pm 0.1$ & $90.2 \pm 0.3$ & $89.8 \pm 0.1$ & $89.9 \pm 0.1$ & $89.9 \pm 0.1$ & $89.8 \pm 0.1$ \\
 \bottomrule
\end{tabular}
\label{tab:delky}
\end{table}

Teplota během měření byla termohygrobarometrem Commeter změřena jako $t_0 = (22.8 \pm 0.4) \, \si{\celsius}$ 

\subsection{Čtyřbodové zapojení}

Nejprve byly odpory drátů měřeny pomocí multimetru KEITHLEY 2010 v čtyřbodovém zapojení. Vnější (proudové) kontakty drátu
byly zapojeny do zdířek multimetru pro měření proudu a vnitřní (napěťové) kontakty naopak do zdířek pro měření napětí. 
Schematicky je toto zapojení vykresleno na obr. \ref{sch:four}.
Měření probíhalo pro rozsah do \SI{10}{\ohm}, přístor ukazoval přímo hodnoty odporu.

Takto naměřené odpory drátů jsou v tabulce \ref{tab:keithley}. Nejistota hodnot byla vzata z nejistoty přístoje 
(viz \ref{pomucky}). Tyto hodnoty byly následně použity jako odhad odporů při měření pomocí Wheastonova a Thomsonova
můstku.

\begin{table}[!hbtp]
\centering
\caption{Odpory drátů z přístroje KEITHLEY 2010 pro čtyřbodové zapojení}
\begin{tabular}{c|cccccc}
 \toprule
 materiál & wolfram & měď & kanthal & železo & mosaz & chromnikl \\
 \midrule
 $R_K (\si{\milli\ohm})$ & $136.9 \pm 0.1$ & $10.59 \pm 0.09$ & $6287.3 \pm 0.5$ & $1475.6 \pm 0.2$ & $220.8 \pm 0.1$ & $1179.27 \pm 0.16$ \\
 \bottomrule
\end{tabular}
\label{tab:keithley}
\end{table}

\subsection{Měření odporu pomocí Wheastonova můstku} \label{wheaston}

Pro měření odporu pomocí Wheastonova můstku byl přístroj MTW nastaven podle pokynů k měření \cite{pokyny},
přičemž poměr odporů $\frac{R_3}{R_4}$ byl zvolen jako $10^{-3}$. Díky tomu byla hodnota měřeného odporu drátu 
(v jednotkách \si{\milli\ohm}) číselně rovna hodnotě odporu $R_2$ (v \si{\ohm}). Ten byl nejprve odhadnut
hodnotou odporu z multimetru KEITHLEY 2010 a následně vyladěn na odporové dekádě tak, aby byl můstek 
v rovnováze, tedy aby galvanometrem netekl proud (tj. jeho výchylka byla nulová).

V tomto zapojení (schéma na obr. \ref{sch:wheast}) jsou v měřeném odporu započteny i odpory přívodních vodičů. 
Jejich odpor byl zjištěn připojením obou vodičů na stejný kontakt (namísto dvou kontaktů na opačných koncích drátů, 
jejichž odpor byl zkoumán) a opět vyladěním hodnoty odporu na odporové dekádě tak, aby můstek byl v rovnováze.
Odpory přívodních vodičů byly tímto způsobem určeny jako $R_0 = (27 \pm 1) \, \si{\milli\ohm}$. Nejistota byla odhadnuta
jako o řád větší než řád kolečka odporové dekády, při jehož otáčení již výchylka galvanometru zůstávala nulová.

O tuto hodnotu byly následně korigovány naměřené odpory. Hodnoty odporů jednotlivých drátů - bez korekce i s ní -
ukazuje tabulka \ref{tab:wheaston}. Nejistoty korigovaných hodnot byly vypočteny jako
\begin{equation}
    \sigma_{R} = \sqrt{\sigma_{R'} + \sigma_{R_0}} \, ,
\end{equation}
výsledná hodnota byla ovšem zaokrouhlena na jednu platnou číslici, protože většina hodnot odporu nebyla určena přesněji, než byl řád odpovídající první platné číslici nejistoty.

\begin{table}[!hbtp]
\centering
\caption{Odpory drátů zjištěné pomocí Wheastonova můstku, bez korekce a korigované o odpor přívodních drátů}
\begin{tabular}{c|cccccc}
 \toprule
 materiál & wolfram & měď & kanthal & železo & mosaz & chromnikl \\
 \midrule
 $R'_W (\si{\milli\ohm})$ & $169.1 \pm 0.1$ & $38.3 \pm 0.1$ & $6311 \pm 1$ & $1508 \pm 1$ & $250 \pm 1$ & $1208 \pm 1$ \\
 $R_W (\si{\milli\ohm})$ & $142 \pm 1$ & $11 \pm 1$ & $6284 \pm 1$ & $1481 \pm 1$ & $223 \pm 1$ & $1181 \pm 1$ \\

 \bottomrule
\end{tabular}
\label{tab:wheaston}
\end{table}




\subsection{Měření odporu pomocí Thomsonova můstku} \label{thomson}

Následně byl přístroj Metra MTW přenastaven do konfigurace pro měření odporu Thomsnovým můstkem. Odpor $R_N$ ze schématu \ref{sch:thomson} měl hodnotu \SI{0.1}{\ohm}, odpory $P$ a $p$ byly nastaveny na stejnou hodnotu \SI{100}{\ohm}. Odporová dekáda je díky přístroji zdvojená a odpovídá odporům $Q$ a $q$ \cite{pokyny}. Díky tomuto nastavení tedy z rovnice \ref{eq:thomson} číselná hodnota na odporové dekádě (v jednotkách \si{\ohm}) odpovídá hodnotě měřeného odporu, avšak v jednotkách \si{\milli\ohm}. Změřené hodnoty odporů jednotlivých drátů, které odpovídají hodnotě odporu na odporové dekádě pro Thomsonův můstek v rovnováze, ukazuje tabulka \ref{tab:thomson} Jejich nejistota byla určena obdobně jako v případě měření s Wheastonovým můstkem.

\begin{table}[!hbtp]
\centering
\caption{Odpory drátů zjištěné pomocí Thomsonova můstku}
\begin{tabular}{c|cccccc}
 \toprule
 materiál & wolfram & měď & kanthal & železo & mosaz & chromnikl \\
 \midrule
 $R_T (\si{\milli\ohm})$ & $137.6 \pm 0.1$ & $10.9 \pm 0.1$ & $6284 \pm 1$ & $1482 \pm 1$ & $221.4 \pm 0.1$ & $1179 \pm 1$ \\
 \bottomrule
\end{tabular}
\label{tab:thomson}
\end{table}

\subsection{Měrný odpor}

Pro všechny tři metody měření odporu byl z rovnic \ref{eq:s} a \ref{eq:rho} určen měrný odpor příslušného drátu. Jeho nejistota byla vypočtena podle \cite{englich} jako
\begin{equation}
    \sigma_\rho = \rho \sqrt{\relunc{R} + \relunc{l} + 4 \relunc{D}} \, .
\end{equation}

Pro srovnání s tabulkovými hodnotami byl z těchto hodnot měrných odporů vypočten aritmetický průměr, jeho nejistota byla vypočtena jako
\begin{equation}
    \sigma_{\bar{\rho}} = \sqrt{\frac{\left(\rho_K - \bar{\rho}\right)^2+ \left(\rho_W - \bar{\rho}\right)^2 + \left(\rho_T - \bar{\rho}\right)^2}{2} + \left(\frac{\sigma_{\rho}}{3}\right)^2} \, ,
\end{equation}
kde $\sigma_\rho$ je průměrná hodnota nejistoty měrného odporu pro jednotlivé metody měření. Protože nejistota měrného odporu je za použití různých metod měření téměř stejná, je s ní zacházeno jako s nejistotou tybu B, tedy se předepsaným způsobem \cite{englich} sčítá s nejistotou aritmetického průměru (vlastně nejistota typu A).


\begin{table}[!htbp]
\centering
\caption{Odpory a měrné odpory jednotlivých drátů pro různé metody měření}
\begin{tabular}{c|cccccc}
 \toprule
 materiál & $R_K\, (\mathrm{m\Omega})$ & $R_W\, (\mathrm{m\Omega})$ & $R_T\, (\mathrm{m\Omega})$ & $\rho_K\, (10^{-8} \, \mathrm{\Omega \cdot m})$ & $\rho_W\, (10^{-8} \, \mathrm{\Omega \cdot m})$ & $\rho_T\, (10^{-8} \, \mathrm{\Omega \cdot m})$ \\
 \midrule
 wolfram & $136.9 \pm 0.1$ & $142.1 \pm 0.1$ & $137.6 \pm 0.1$ & $5.64 \pm 0.18$ & $5.86 \pm 0.19$ & $5.67 \pm 0.18$ \\
 měď & $10.59 \pm 0.09$ & $11.3 \pm 0.1$ & $10.9 \pm 0.1$ & $1.67 \pm 0.03$ & $1.78 \pm 0.03$ & $1.72 \pm 0.03$ \\
 kanthal & $6287.3 \pm 0.5$ & $6284 \pm 1$ & $6284 \pm 1$ & $138 \pm 6$ & $138 \pm 6$ & $138 \pm 6$ \\
 železo & $1475.6 \pm 0.2$ & $1481 \pm 1$ & $1482 \pm 1$ & $21.7 \pm 1.1$ & $21.7 \pm 1.1$ & $21.8 \pm 1.1$ \\
 mosaz & $220.8 \pm 0.1$ & $223 \pm 1$ & $221.4 \pm 0.1$ & $6.9 \pm 0.3$ & $6.9 \pm 0.3$ & $6.9 \pm 0.3$ \\
 chromnikl & $1179.27 \pm 0.16$ & $1181 \pm 1$ & $1179 \pm 1$ & $104 \pm 3$ & $104 \pm 3$ & $104 \pm 3$ \\
 \bottomrule
\end{tabular}
\label{tab:rho}
\end{table}

\begin{table}[!hbtp]
\centering
\caption{Střední hodnota měrného odporu pro jednotlivé materiály}
\begin{tabular}{c|cccccc}
 \toprule
 materiál & wolfram & měď & kanthal & železo & mosaz & chromnikl \\
 \midrule
 $\bar{\rho} \, (10^{-8} \, \mathrm{\Omega \cdot m})$ & $5.7 \pm 0.2$ & $1.72 \pm 0.06$ & $138 \pm 6$ & $21.7 \pm 1.1$ & $6.9 \pm 0.3$ & $104 \pm 3$ \\
 \bottomrule
\end{tabular}
\label{tab:meanrho}
\end{table}



% ============================ DISKUZE ===============================
\section{Diskuze}

Odpory drátů z různých materiálů byly měřeny pomocí tří odlišných metod: multimetru Keithley 2010, Wheastonova můstku a Thomsonova můstku. Nejvíce odlišné byly hodnoty odporu z Wheastonova můstku, protože byly systematicky větší o odpory přívodních drátů a odpory svorek. Při měření s multimetrem Keithley a Thomsonovým můstkem se totiž využívalo čtyřbodové zapojení rezistoru (schéma \ref{sch:four}), při kterém se napětí měří mezi napěťovými kontakty a odpor přívodních drátů se tak neuplatní. Odpor přívodních drátů byl u Wheastonova můstku změřen jako $R_0 = (27 \pm 1) \, \si{\milli\ohm}$, což je téměř třikrát větší hodnota než byl skutečný odpor měděného drátu. A i pro ostatní dráty se již tato hodnota pohybuje v řádu jednotek až desítek procent měřené hodnoty, nelze jej tedy zanedbat a bylo jej nutné od výsledných hodnot odečíst. V tomto dodatečném odporu je obsažen odpor přívodních drátů i odpor svorek, na které byly dráty připojeny. Jeho hodnota však nebyla měřena pro všechny svorky, je tedy možné, že se odpor různých svorek liší, a tedy odpor změřený pomocí Wheastonova můstku by byl stále zatížen systematickou chybou. Kromě wolframového drátu však byly hodnoty odporu ze~všech metod měření v rámci nejistoty velmi podobné. Pokud se tedy odpory jednotlivých svorek liší, zřejmě není tento rozdíl příliš velký a nemá významný vliv na změřený odpor drátu.

Měření Wheastonovým můstkem bylo specifické i tím, že se při dlouhém zapojení drátu do obvodu začal drát ohřívat, čímž jeho odpor lehce rostl. Bylo pozorováno, že výchylka galvanometru - nejprve nulová - se začala s rostoucí dobou zapojení drátu do obvodu zvětšovat. Proto bylo nutné měřit galvanometrem (spojit obvod a nechat jím téct proud) vždy jen krátký okamžik. U měření s Thomsonovým můstkem tekl proud obvodem se zkoumaným drátem neustále a změna výchylky galvanometru při delším sepnutí nebyla pozorována. Díky nízké hodnotě proudu se však drát a rezistory v obvodu pravděpodobně nezahřívaly tolik jako při měření s Wheastonovým můstkem. I tak ale může takovéto ohřátí vodiče procházejícím proudem systematicky zvýšit hodnotu zkoumaného odporu.

Do výsledné hodnoty měrného odporu vnášela (u většiny materiálů) největší nejistotu hodnota průměru drátu. Jednak se kvůli přesnosti měřidla nedá určit dostatečně přesně, jednak se ve vztahu pro výpočet měrného odporu objevuje ve druhé mocnině a má tak ze všech veličin největší váhu. Mikrometrem byly průměry drátů určeny s přesností v řádu jednotek procent, největší relativní přesnosti bylo dosaženo u měděného drátu (0,8\,\%) a nejmenší u železného (2,4\,\%) díky jejich velkému resp. malému průměru. Měrný odpor měděného drátu byl nejvíce zatížen nejistotou celkového odporu, a to kvůli jeho velmi nízké hodnotě. Měděný drát nebyl mezi svorkami napnut úplně dokonale, jeho délka a i odhadovaná nejistota byly větší než u ostatních drátů. V~relativním měřítku však u měděného drátu převážila nejistota odporu: asi 0,9\,\%. Relativní přesnosti celkových odporů drátů se lehce lišily v závislosti na použité metodě měření, obecně bylo nejpřesnější měření multimetrem Keithley 2010.

Měření probíhalo za teploty $t_0 = (22.8 \pm 0.4) \, \si{\celsius}$. Vzhledem k tomu, že odpor drátů závisí na jeho teplotě, bylo nutné před srovnáním experimentální a tabulkové hodnoty měrného odporu provést její přepočet na stejnou teplotu. Experimentálně určená hodnota měrného odporu wolframu je $\bar{\rho} = (5.7 \pm 0.2) \cdot 10^{-8} \, \mathrm{\Omega \cdot m}$. Tabulková hodnota při teplotě \SI{0}{\celsius} je podle \cite{mikulcak} $\rho' = 5.3 \cdot 10^{-8} \, \mathrm{\Omega \cdot m}$ s teplotním součinitelem odporu $\alpha = 4.4 \cdot 10^{-3} \, \mathrm{K^{-1}}$. Za předpokladu lineární závislosti měrného odporu na teplotě vychází tabulková hodnota pro $t = 22.8 \,^\circ \mathrm{C}$ $\rho'' = 5.8 \cdot 10^{-8} \, \mathrm{\Omega \cdot m}$. Hodnoty se tedy v rámci nejistoty shodují. 

Pro měď je změřená hodnota $\bar{\rho} = (1.72 \pm 0.06) \cdot 10^{-8} \, \mathrm{\Omega \cdot m}$ a tabulková pro \SI{0}{\celsius} je podle \cite{mikulcak} $\rho' = 1.8 \cdot 10^{-8} \, \mathrm{\Omega \cdot m}$. S teplotním součinitelem odporu $\alpha = 4.0 \cdot 10^{-3} \, \mathrm{K^{-1}}$ je přepočtená hodnota $\rho'' = 2.0 \cdot 10^{-8} \, \mathrm{\Omega \cdot m}$. Měrný odpor bude jistě záviset na čistotě materiálu, u mědi jakožto velmi dobrého vodiče, jež je vyznačuje nízkým měrným odporem, obzvláště. Vzhledem k tomu, že řádově si hodnoty odpovídají, lze experimentálně určenou hodnotu považovat za správnou.

Experimentální hodnota měrného odporu kanthalu je $\bar{\rho} = (138 \pm 6) \cdot 10^{-8} \, \mathrm{\Omega \cdot m}$ a tabulková při \SI{0}{\celsius} je podle \cite{mikulcak} $\rho' = 140 \cdot 10^{-8} \, \mathrm{\Omega \cdot m}$ s teplotním součinitelem odporu $\alpha = 0.06 \cdot 10^{-3} \, \mathrm{K^{-1}}$, tedy hodnota pro teplotu \SI{22.8}{\celsius} bude téměř stejná. V rámci nejistoty (která je u kanthalu daná především nejistotou průměru drátu) se obě hodnoty shodují. Zkoumaná slitina má tedy pravděpodobně velmi podobné složení jako slitina z tabulek.

Měrný odpor železa byl určen jako $\bar{\rho} = (21.7 \pm 1.1) \cdot 10^{-8} \, \mathrm{\Omega \cdot m}$. Tabulková hodnota je podle \cite{zelezo} \\ $\rho' = 9.8 \cdot 10^{-8} \, \mathrm{\Omega \cdot m}$ při \SI{20}{\celsius} a teplotní součinitel odporu je $\alpha = 6 \cdot 10^{-3} \, \mathrm{K^{-1}}$. Po přepočtu na teplotu během měření vychází hodnota $\rho'' = 10.0 \cdot 10^{-8} \, \mathrm{\Omega \cdot m}$. I přepočtená hodnota je tedy přibližně poloviční oproti té experimentální. Mechanické a elektrické vlastnosti železa však velmi významným způsobem závisejí nejen na čistotě, ale také na mechanickém a tepelném zpracování a postupu výroby. Vzhledem k tomu, že měrný odpor byl měřen třemi různými způsoby a výsledek tak pravděpodobně není ovlivněn hrubou chybou, můžeme soudit, že železný drát použitý při měření byl vyroben odlišným způsobem než výrobek studovaný v tabulkách \cite{zelezo}.

Měrný odpor mosazi byl zjištěn jako $\bar{\rho} = (6.9 \pm 0.3) \cdot 10^{-8} \, \mathrm{\Omega \cdot m}$. Tabulková hodnota při \SI{0}{\celsius} z \cite{mikulcak} je $\rho' = 8 \cdot 10^{-8} \, \mathrm{\Omega \cdot m}$ a teplotní součinitel odporu je $\alpha = 1.5 \cdot 10^{-3} \, \mathrm{K^{-1}}$. Přepočtená hodnota na \SI{22.8}{\celsius} je $\rho'' = 8.3 \cdot 10^{-8} \, \mathrm{\Omega \cdot m}$. Experimentální a tabulková hodnota se liší. To by mohlo být dané různým zastoupením prvků ve zkoumané slitině a té z tabulek, případně odlišným procesem výroby. Vlastnosti slitin vyrobené různými postupy se totiž mohou i docela výrazně lišit. Opět jsou však hodnoty řádově podobné.

Pro chromnikl byl měrný odpor určen jako $\bar{\rho} = (104 \pm 3) \cdot 10^{-8} \, \mathrm{\Omega \cdot m}$. Tabulky \cite{mikulcak} udávají hodnotu \\ $\rho' = 110 \cdot 10^{-8} \, \mathrm{\Omega \cdot m}$ při \SI{0}{\celsius} a teplotní součitnitel odporu $\alpha = 0.18 \cdot 10^{-3} \, \mathrm{K^{-1}}$. Se změnou o \SI{22.8}{\celsius} se tedy hodnota odporu příliš nezmění. Tabulková hodnota se od té experimentální liší od 2 směrodatné odchylky. Po uvážení vlivu složení slitiny na její mechanické a elektrické vlastnosti lze i hodnotu měrného elektrického odporu chromniklu považovat za správnou.


% ============================= ZÁVĚR ================================
\section{Závěr}

Pro dráty z různých materiálů byly různými metodami (multimetrem Keithley 2010 ve čtyřbodovém zapojení, Wheastonovým můstkem a Thomsonovým můstkem) měřeny odpory a měrné odpory.

Pro wolfram byly změřeny hodnoty odporu
\begin{equation*}
    R_K = (136.9 \pm 0.1) \, \si{\milli\ohm}
\end{equation*}
\begin{equation*}
    R_W = (142.1 \pm 0.1) \, \si{\milli\ohm}
\end{equation*}
\begin{equation*}
    R_T = (137.6 \pm 0.1) \, \si{\milli\ohm}
\end{equation*}
a vypočteny hodnoty měrného odporu
\begin{equation*}
    \rho_K = (5.64 \pm 0.18) \cdot 10^{-8} \, \mathrm{\Omega \cdot m}
\end{equation*}
\begin{equation*}
    \rho_W = (5.86 \pm 0.19) \cdot 10^{-8} \, \mathrm{\Omega \cdot m}
\end{equation*}
\begin{equation*}
    \rho_T = (5.67 \pm 0.18) \cdot 10^{-8} \, \mathrm{\Omega \cdot m}
\end{equation*}
se střední hodnotou
\begin{equation*}
    \bar{\rho} = (5.7 \pm 0.2) \cdot 10^{-8} \, \mathrm{\Omega \cdot m} \, .
\end{equation*}

Pro měď byly změřeny hodnoty odporu
\begin{equation*}
    R_K = (10.59 \pm 0.09) \, \si{\milli\ohm}
\end{equation*}
\begin{equation*}
    R_W = (11.3 \pm 0.1) \, \si{\milli\ohm}
\end{equation*}
\begin{equation*}
    R_T = (10.9 \pm 0.1) \, \si{\milli\ohm}
\end{equation*}
a vypočteny hodnoty měrného odporu
\begin{equation*}
    \rho_K = (1.67 \pm 0.03) \cdot 10^{-8} \, \mathrm{\Omega \cdot m}
\end{equation*}
\begin{equation*}
    \rho_W = (1.78 \pm 0.03) \cdot 10^{-8} \, \mathrm{\Omega \cdot m}
\end{equation*}
\begin{equation*}
    \rho_T = (1.72 \pm 0.03) \cdot 10^{-8} \, \mathrm{\Omega \cdot m}
\end{equation*}
se střední hodnotou
\begin{equation*}
    \bar{\rho} = (1.72 \pm 0.06) \cdot 10^{-8} \, \mathrm{\Omega \cdot m} \, .
\end{equation*}

Pro kanthal byly změřeny hodnoty odporu
\begin{equation*}
    R_K = (6287.3 \pm 0.5) \, \si{\milli\ohm}
\end{equation*}
\begin{equation*}
    R_W = (6284 \pm 1) \, \si{\milli\ohm}
\end{equation*}
\begin{equation*}
    R_T = (6284 \pm 1) \, \si{\milli\ohm}
\end{equation*}
a vypočteny hodnoty měrného odporu
\begin{equation*}
    \rho_K = (138 \pm 6) \cdot 10^{-8} \, \mathrm{\Omega \cdot m}
\end{equation*}
\begin{equation*}
    \rho_W = (138 \pm 6) \cdot 10^{-8} \, \mathrm{\Omega \cdot m}
\end{equation*}
\begin{equation*}
    \rho_T = (138 \pm 6) \cdot 10^{-8} \, \mathrm{\Omega \cdot m}
\end{equation*}
se střední hodnotou
\begin{equation*}
    \bar{\rho} = (138 \pm 6) \cdot 10^{-8} \, \mathrm{\Omega \cdot m} \, .
\end{equation*}

Pro železo byly změřeny hodnoty odporu
\begin{equation*}
    R_K = (1475.6 \pm 0.2) \, \si{\milli\ohm}
\end{equation*}
\begin{equation*}
    R_W = (1481 \pm 1) \, \si{\milli\ohm}
\end{equation*}
\begin{equation*}
    R_T = (1482 \pm 1) \, \si{\milli\ohm}
\end{equation*}
a vypočteny hodnoty měrného odporu
\begin{equation*}
    \rho_K = (21.7 \pm 1.1) \cdot 10^{-8} \, \mathrm{\Omega \cdot m}
\end{equation*}
\begin{equation*}
    \rho_W = (21.7 \pm 1.1) \cdot 10^{-8} \, \mathrm{\Omega \cdot m}
\end{equation*}
\begin{equation*}
    \rho_T = (21.8 \pm 1.1) \cdot 10^{-8} \, \mathrm{\Omega \cdot m}
\end{equation*}
se střední hodnotou
\begin{equation*}
    \bar{\rho} = (21.7 \pm 1.1) \cdot 10^{-8} \, \mathrm{\Omega \cdot m} \, .
\end{equation*}

Pro mosaz byly změřeny hodnoty odporu
\begin{equation*}
    R_K = (220.8 \pm 0.1) \, \si{\milli\ohm}
\end{equation*}
\begin{equation*}
    R_W = (223 \pm 1) \, \si{\milli\ohm}
\end{equation*}
\begin{equation*}
    R_T = (221.4 \pm 0.1) \, \si{\milli\ohm}
\end{equation*}
a vypočteny hodnoty měrného odporu
\begin{equation*}
    \rho_K = (6.9 \pm 0.3) \cdot 10^{-8} \, \mathrm{\Omega \cdot m}
\end{equation*}
\begin{equation*}
    \rho_W = (6.9 \pm 0.3) \cdot 10^{-8} \, \mathrm{\Omega \cdot m}
\end{equation*}
\begin{equation*}
    \rho_T = (6.9 \pm 0.3) \cdot 10^{-8} \, \mathrm{\Omega \cdot m}
\end{equation*}
se střední hodnotou
\begin{equation*}
    \bar{\rho} = (6.9 \pm 0.3) \cdot 10^{-8} \, \mathrm{\Omega \cdot m} \, .
\end{equation*}

Pro chromnikl byly změřeny hodnoty odporu
\begin{equation*}
    R_K = (1179.27 \pm 0.16) \, \si{\milli\ohm}
\end{equation*}
\begin{equation*}
    R_W = (1181 \pm 1) \, \si{\milli\ohm}
\end{equation*}
\begin{equation*}
    R_T = (1179 \pm 1) \, \si{\milli\ohm}
\end{equation*}
a vypočteny hodnoty měrného odporu
\begin{equation*}
    \rho_K = (104 \pm 3) \cdot 10^{-8} \, \mathrm{\Omega \cdot m}
\end{equation*}
\begin{equation*}
    \rho_W = (104 \pm 3) \cdot 10^{-8} \, \mathrm{\Omega \cdot m}
\end{equation*}
\begin{equation*}
    \rho_T = (104 \pm 3) \cdot 10^{-8} \, \mathrm{\Omega \cdot m}
\end{equation*}
se střední hodnotou
\begin{equation*}
    \bar{\rho} = (104 \pm 3) \cdot 10^{-8} \, \mathrm{\Omega \cdot m} \, .
\end{equation*}

% =========================== LITERATURA =============================
\begin{thebibliography}{9}

\bibitem{studijni} Kolektiv ZFP KVOF MFF UK: Měření malých odporů [online]. [cit. 24.11.2025], dostupné z~https://physics.mff.cuni.cz/vyuka/zfp/\_media/zadani/texty/txt\_204.pdf

\bibitem{pokyny} Kolektiv ZFP KVOF MFF UK: Měření malých odporů - pokyny k měření [online]. [cit. 24.11.2025], dostupné
z~https://physics.mff.cuni.cz/vyuka/zfp/\_media/zadani/pokyny/mereni\_204.pdf

\bibitem{englich} J. Englich: Úvod do praktické fyziky I: Zpracování výsledků měření. 1. vyd. Praha: Matfyzpress, 2006

\bibitem{mikulcak} Mikulčák a kolektiv: Matematické, fyzikální a chemické tabulky pro střední školy: relativní permitivita, Praha: Státní pedagogické nakladatelství, n. p., 1988

\bibitem{zelezo} J. Bureš: conVERTER - převody jednotek: měrný odpor [online]. [cit. 24.11.2025], dostupné z~https://www.converter.cz/tabulky/merny-odpor.htm

\end{thebibliography}

\end{document}